\section{선별 문제 풀이 II}
\label{sec:equation-groups-solutions-b}

\subsection*{문제 \thechapter.16: 기약 삼차식과 판별식}

\begin{solution}
$f=X^3+aX+b$가 기약이므로 갈루아군 $G$는 세 근 위에서 추이적으로 작용한다. $S_3$의 추이적 부분군은 $A_3$와 $S_3$뿐이다.

근들을 $\alpha_1,\alpha_2,\alpha_3$라 하고
\[
  \delta=(\alpha_1-\alpha_2)(\alpha_1-\alpha_3)(\alpha_2-\alpha_3)
\]
라 두자. 순열 $\sigma$는
\[
  \sigma(\delta)=\sgn(\sigma)\delta
\]
를 만족한다. 따라서
\[
  G\subseteq A_3
  \quad\Longleftrightarrow\quad
  \delta\text{가 }G\text{ 전체에 의해 고정됨}
  \quad\Longleftrightarrow\quad
  \delta\in K.
\]
한편
\[
  \delta^2=\Delta_f=-4a^3-27b^2.
\]
그러므로 $\Delta_f$가 $K$의 제곱일 때 정확히 $\delta\in K$이고, 이 경우 추이성 때문에 $G=A_3\cong C_3$이다. 제곱이 아니면 $G$는 홀순열을 포함하므로 $G=S_3$이다.
\end{solution}

\subsection*{문제 \thechapter.18: 변수 순열과 아르틴의 정리}

\begin{solution}
$S_n$은
\[
  \pi(T_i)=T_{\pi(i)}
\]
로 $L=k(T_1,\dots,T_n)$에 작용한다. 이 작용이 자명하다면 모든 $T_i$가 고정되므로 $\pi(i)=i$가 모든 $i$에 대해 성립하고 $\pi=1$이다. 따라서 작용은 충실하고 $S_n$은 $\Aut_k(L)$의 위수 $n!$인 부분군이다.

고정체를 $K=L^{S_n}$이라 하자. 아르틴의 고정체 정리를 적용하면
\[
  L/K\text{는 유한 갈루아},
  \qquad
  \Gal(L/K)=S_n,
  \qquad
  [L:K]=|S_n|=n!.
\]
\end{solution}

\subsection*{문제 \thechapter.21: 안정자와 단순 중간체}

\begin{solution}
$\sigma\in G$가 $\alpha$를 고정하면 $K$의 원소도 고정하므로 모든 유리식 $h(\alpha)$를 고정한다. 따라서 $K(\alpha)$를 점별로 고정한다. 역으로 $K(\alpha)$를 점별로 고정하는 자동동형은 당연히 $\alpha$를 고정한다. 그러므로
\[
  \Stab_G(\alpha)=\Gal(L/K(\alpha)).
\]
유한 갈루아 대응의 차수 공식에서
\[
  [K(\alpha):K]
  =[G:\Gal(L/K(\alpha))]
  =[G:\Stab_G(\alpha)].
\]
이는 궤도-안정자 공식과 함께
\[
  [K(\alpha):K]=|G\cdot\alpha|
\]
를 준다.
\end{solution}

\subsection*{문제 \thechapter.23: 음의 판별식을 가진 기약 삼차식}

\begin{solution}
$\Q$의 제곱은 음수가 될 수 없으므로 $\Delta_f<0$이면 $\Delta_f$는 $\Q$의 제곱이 아니다. 기약 삼차식의 판별식 판정에 의해
\[
  \Gal(f/\Q)\cong S_3.
\]

근의 실수성도 같은 현상을 보여준다. 세 근이 모두 실수이면
\[
  \delta=(\alpha_1-\alpha_2)(\alpha_1-\alpha_3)(\alpha_2-\alpha_3)
\]
가 실수이므로 $\Delta_f=\delta^2>0$이다. 따라서 $\Delta_f<0$인 분리 삼차식은 실근 하나와 서로 켤레인 비실근 두 개를 갖는다. 복소켤레는 두 비실근을 교환하는 전치, 즉 홀순열을 갈루아군 안에 제공한다. 추이적 부분군이 $3$-순환과 전치를 모두 포함하므로 $S_3$가 된다.
\end{solution}

\subsection*{문제 \thechapter.25: 일반다항식의 갈루아군}

\begin{solution}
독립변수 $T_1,\dots,T_n$과
\[
  L=k(T_1,\dots,T_n)
\]
을 잡는다. $S_n$이 변수들을 순열하도록 하고 기본대칭식들을 $s_1,\dots,s_n$이라 하자. 아르틴 고정체 정리와 대칭 유리함수 정리에 의해
\[
  L^{S_n}=k(s_1,\dots,s_n),
  \qquad
  \Gal(L/k(s_1,\dots,s_n))=S_n.
\]
또한
\[
  \prod_{i=1}^n(X-T_i)
  =X^n-s_1X^{n-1}+\cdots+(-1)^ns_n.
\]

$s_1,\dots,s_n$은 대수적으로 독립이므로 독립변수 $S_1,\dots,S_n$과 부호를 맞추어
\[
  S_j\longmapsto(-1)^j s_j
\]
인 체동형을 정의할 수 있다. 이 동형 아래 일반다항식
\[
  P(X)=X^n+S_1X^{n-1}+\cdots+S_n
\]
은 위 곱으로 옮겨진다. 따라서 그 분해체와 갈루아군도 동형이고
\[
  \Gal(P/k(S_1,\dots,S_n))\cong S_n.
\]
$S_n$의 근 작용은 추이적이므로 $P$는 기약이고, 근 $T_i$들이 서로 다른 독립변수이므로 분리이다.
\end{solution}

\subsection*{문제 \thechapter.27: $X^3-X+1$의 중간체}

\begin{solution}
가능한 유리근 $\pm1$을 대입하면 영이 아니므로 $f=X^3-X+1$은 기약이다. 판별식은
\[
  \Delta_f=4-27=-23
\]
이고 $\Q$의 제곱이 아니므로 분해체 $L$에 대해
\[
  G=\Gal(L/\Q)\cong S_3,
  \qquad [L:\Q]=6.
\]

$S_3$의 부분군은 자명부분군, 세 개의 위수 $2$ 부분군, $A_3$, 전체 군이다. 따라서 중간체는 다음과 같다.

\begin{itemize}
  \item $L$ 자체: 차수 $6$, 정규.
  \item 세 근 $\alpha_1,\alpha_2,\alpha_3$에 대응하는 세 삼차체 $\Q(\alpha_i)$: 각각 차수 $3$, 비정규.
  \item $A_3$의 고정체: 차수 $2$, 정규.
  \item $\Q$: 차수 $1$.
\end{itemize}

교대곱 $\delta$는 $A_3$가 고정하고 $\delta^2=-23$이므로 유일한 이차 중간체는
\[
  \Q(\delta)=\Q(\sqrt{-23}).
\]
\end{solution}

\subsection*{문제 \thechapter.29: 열두 번째 원분체의 중간체}

\begin{solution}
$L=\Q(\zeta_{12})$이고
\[
  G\cong(\Z/12\Z)^\times
  =\{1,5,7,11\}
  \cong V_4.
\]
$a\in\{1,5,7,11\}$에 대해 $\sigma_a(\zeta)=\zeta^a$라 쓰자. 세 위수 $2$ 부분군과 고정체는 다음과 같다.

\[
\begin{array}{c|c|c}
H & \text{고정되는 생성원} & L^H\\
\hline
\langle\sigma_5\rangle
& i=\zeta^3
& \Q(i)\\
\langle\sigma_7\rangle
& \zeta^4
& \Q(\zeta^4)=\Q(\sqrt{-3})\\
\langle\sigma_{11}\rangle
& \zeta+\zeta^{-1}=\sqrt3
& \Q(\sqrt3)
\end{array}
\]
각 고정체는 갈루아 대응에서 차수 $2$이므로 표의 포함은 등호이다. 자명부분군과 전체 군은 각각 $L$과 $\Q$에 대응한다. $G$가 아벨군이므로 모든 부분군이 정규이고, 따라서 모든 중간체는 $\Q$ 위 갈루아이다.
\end{solution}

\subsection*{문제 \thechapter.31: 인수 차수와 궤도 크기}

\begin{solution}
$f$가 분리이고
\[
  f=c f_1\cdots f_r
\]
가 서로 다른 기약인수들의 곱이라 하자. 각 $f_i$의 근 집합은 $K$-계수 조건 때문에 갈루아군에 의해 보존된다. $f_i$가 기약이므로 그 근들 위의 작용은 추이적이다. 따라서 $f_i$의 근 집합은 정확히 하나의 궤도이고 그 크기는 $\deg f_i=d_i$이다.

역으로 근 집합의 궤도 하나를 택하고 한 근 $\alpha$를 고르면, 그 궤도는 $\alpha$의 최소다항식의 모든 근으로 이루어진다. 따라서 궤도 크기는 최소다항식 차수이며, 궤도들로부터 기약인수 차수들을 복원할 수 있다.

예제
\[
  (X^2-2)(X^3-3)(X-1)
\]
에서 세 인수는 $\Q$ 위에서 기약이고 분리이다. 따라서 궤도 크기는
\[
  2,\quad3,\quad1
\]
이다. 구체적으로
\[
  \{\sqrt2,-\sqrt2\},
  \quad
  \{\sqrt[3]{3},\omega\sqrt[3]{3},\omega^2\sqrt[3]{3}\},
  \quad
  \{1\}
\]
이 세 궤도이다.
\end{solution}
